\subsection{Euclid's Five Postulates}

The following are the five postulates used in the first book of Euclid's
\emph{Elements}. They specify the elementary geometric operations and relations
from which the later constructions begin. The diagrams illustrate their meaning;
they are not proofs.

\begin{corebox}[Postulate I: joining two points]
A straight segment can be drawn from any point to any other point.
\end{corebox}

\begin{center}
\begin{tikzpicture}[scale=0.9]
    \coordinate (A) at (-2.4,-0.6);
    \coordinate (B) at (2.1,0.8);
    \draw[subset] (A) -- (B);
    \fill[geometricSubsetColor] (A) circle (2pt) node[below] {$A$};
    \fill[geometricSubsetColor] (B) circle (2pt) node[above] {$B$};
    \node at (0,-1.55) {a segment is drawn from $A$ to $B$};
\end{tikzpicture}
\end{center}

\begin{corebox}[Postulate II: extending a segment]
A finite straight segment can be extended continuously in the same straight
line.
\end{corebox}

\begin{center}
\begin{tikzpicture}[scale=0.9]
    \coordinate (A) at (-2.2,0);
    \coordinate (B) at (1.2,0);
    \coordinate (C) at (3.1,0);
    \draw[subset] (A) -- (B);
    \draw[helper,->] (B) -- (C);
    \fill[geometricSubsetColor] (A) circle (2pt) node[below] {$A$};
    \fill[geometricSubsetColor] (B) circle (2pt) node[below] {$B$};
    \node[below] at (C) {$C$};
    \node at (0,-1.15) {$AB$ is extended beyond $B$};
\end{tikzpicture}
\end{center}

\begin{corebox}[Postulate III: drawing a circle]
A circle can be drawn with any chosen centre and any chosen radius.
\end{corebox}

\begin{center}
\begin{tikzpicture}[scale=0.9]
    \coordinate (O) at (0,0);
    \coordinate (A) at (2.2,0);
    \draw[helper] (O) circle (2.2);
    \draw[subset] (O) -- (A) node[midway,below] {$r$};
    \fill[geometricSubsetColor] (O) circle (2pt) node[below left] {$O$};
    \fill[geometricSubsetColor] (A) circle (2pt) node[right] {$A$};
    \node at (0,-2.75) {centre $O$ and radius $r$ determine a circle};
\end{tikzpicture}
\end{center}

\begin{interpretationbox}[Why the circle postulate matters]
This postulate permits a length to be copied geometrically without first
assigning a number to it. The compass transfers a geometric magnitude directly.
\end{interpretationbox}

\begin{corebox}[Postulate IV: equality of right angles]
All right angles are congruent.
\end{corebox}

\begin{center}
\begin{tikzpicture}[scale=0.85]
    \draw[subset] (-3.2,0) -- (-0.2,0);
    \draw[subset] (-1.7,-1.3) -- (-1.7,1.6);
    \draw (-1.42,0) -- (-1.42,0.28) -- (-1.7,0.28);

    \draw[subset] (0.8,-0.8) -- (3.5,0.7);
    \draw[subset] (2.9,-1.25) -- (1.4,1.45);
    \draw (2.24,0.02) -- (2.10,0.27) -- (1.85,0.13);

    \node at (0,-2.0) {every right angle represents the same geometric magnitude};
\end{tikzpicture}
\end{center}

\begin{remarkbox}
This postulate does not say that a perpendicular may simply be drawn without
construction. It says that once right angles have been constructed, they are all
congruent.
\end{remarkbox}

\begin{corebox}[Postulate V: the parallel postulate]
If a transversal meets two straight lines and the interior angles on one side
sum to less than two right angles, then the two lines, when extended, meet on
that side.
\end{corebox}

\begin{center}
\begin{tikzpicture}[scale=0.85]
    \draw[subset] (-3.4,1.15) -- (3.5,0.15) node[right] {$\ell$};
    \draw[subset] (-3.4,-1.2) -- (3.5,-0.35) node[right] {$m$};
    \draw[helper] (-1.0,-2.0) -- (0.6,2.0) node[above] {$t$};
    \node at (0,-2.45) {the angle condition determines on which side $\ell$ and $m$ meet};
\end{tikzpicture}
\end{center}

\begin{interpretationbox}[Modern working form]
A familiar equivalent is Playfair's formulation: through a point outside a
given line there passes exactly one parallel to that line. We shall use that
form later, after the relation between the two formulations has been made clear.
\end{interpretationbox}

\begin{summarybox}[What the postulates provide]
The first three postulates provide elementary constructions, the fourth makes
right angles comparable, and the fifth controls parallelism. Together they
specify the geometric world available before coordinates and algebra are
introduced.
\end{summarybox}
